\scanpage{299}
- \mspage{3} -
Tribu n$^\circ$ 43

(avec les ``dans le cas séparable c'est déjà connu''. Le sorite du no 2 est assez agréable pour les extensions régulières ; on pourrait donc commencer par celles-ci, et profiter de ce qu'on saura alors pour traiter les extensions primaires.

Chevalley a vendu la démonstration suivante du \uline{lemme de Zariski}. Avec les notations du texte, on suppose $K,E,K(X),E(X)$ plongés dans un domaine universel U (algt. clos) ; soit $u=a(X)/b(X)$ dans $E(X)$, algébrique sur $K(X)$ ; il y a une infinité d'automorphismes s de U sur $K(X)$ ; les éléments $s(u)$ sont en nombre fini ; donc l'ensemble des s tels que $s(u)=u$ est un sous groupe d'indice fini du groupe des $K(X)$-automorphismes de U ; on en déduit que $a(X)=a^{\unsure{s|S}}(X)$ et $b(X)=b^{\unsure{s|S}}(X)$ pour les s de ce sous-groupe, si on a pris la précaution de supposer les polynômes a et b étrangers ; donc a et b ont leurs coefficients algébriques sur K, d'où etcétéra.

Il faut aussi montrer que si K est algébriquement fermé (ou quasi) dans E, et si K' est une extension algébrique finie de K, alors K' est algébriquement fermé (ou quasi) dans E(K'). Voici une démonstration valable en caract. 0, et qu'on pourrait essayer d'arranger : soit K'' la fermeture algébrique de K dans E(K') ; on a $K'\subset K''$ ; or $[E(K'):E]=$ $=[K':K]$ par disjonction linéaire ; de même $[E(K''):E]=[K'':K]$ comme $E(K'')=E(K')$, on en déduit $K''=K'$. C.Q.F.D.

Au n$^\circ$ 4 il faudra expliciter le cas des idéaux principaux (polynômes irréductibles). La remarque d'équidimensionalité ne semble pas plus qu'un exercice. Le fourbi de la p.10 a l'air vrai avec l'hypothèse plus faible que tout élément de S qui est diviseur de 0 dans R l'est dans S (noter que ceci équilibre l'énoncé, qui dit que la propriété se transmet par produit tensoriel).

\mspage{298}
